% GENERATED FILE. Edit canonical lessons, then run npm run generate:books.
% canonical-source-hash: fnv1a64:874336fcc2716f33
% canonical-lessons: SA-C375-vipattih, SA-C375-apad, SA-C375-vinasah, SA-C375-anavrstih, SA-C375-pratikaroti

\chapter{In Hard Times}
\label{ch:sa-in-hard-times}
\hlchaptermodality{375}

\begin{chapteropening}
\textbf{By the end of this chapter:} \emph{I can say what harm will follow if nothing is done.}

Complete the last lesson of chapter 375: I can say what harm will follow if nothing is done.
\end{chapteropening}

\section[\texorpdfstring{\sk{विपत्तिः} (vipattiḥ)}{विपत्तिः (vipattiḥ)}]{\sk{विपत्तिः} (vipattiḥ) — adversity; a calamity}

\label{lesson:SA-C375-vipattih}



\begin{quote}
Before the new one: say the Sanskrit for a mishap, then the Sanskrit for ceases.
\end{quote}

\subsection*{You'll want to know: \sk{विपत्तिः}}
\textbf{\sk{विपत्तिः}} — \emph{vipattiḥ} — \textquotedblleft{}adversity; a calamity\textquotedblright{}.

\begin{quote}
\sk{यदि} \sk{विपत्तिः} \sk{आगच्छति}, \sk{तर्हि} \sk{मित्राणि} \sk{साहाय्यं} \sk{करिष्यन्ति।}
\end{quote}

\begin{quote}
\emph{yadi vipattiḥ āgacchati, tarhi mitrāṇi sāhāyyaṃ kariṣyanti.}
\end{quote}

\begin{quote}
If adversity comes, friends will help.
\end{quote}

A feminine noun.

\begin{cousinweb}
\textbf{\sk{विपत्तिः}} (\emph{vipattiḥ}) is built on \emph{vi}, \textquotedblleft{}apart\textquotedblright{}, and the root √pad, \textquotedblleft{}to go, to fall\textquotedblright{}; MW gives \emph{vi-patti} \textquotedblleft{}going wrongly, adversity, misfortune, failure, disaster\textquotedblright{} (MW).
\end{cousinweb}

\begin{grammarlens}[title={What you've built}]
One more word for saying what happens if.
\end{grammarlens}

\subsection*{Your turn}
\begin{itemize}
\raggedright
  \item \emph{Say it:} \emph{vipattiḥ}
  \item \emph{Say it:} \emph{vipattiḥ}, once more
  \item \emph{Say it:} \emph{yadi vipattiḥ āgacchati, tarhi mitrāṇi sāhāyyaṃ kariṣyanti.} — If adversity comes, friends will help.
  \item \emph{Recall:} say the Sanskrit for a mishap, then the Sanskrit for ceases, then say \emph{vipattiḥ} again
\end{itemize}

\begin{tcolorbox}[breakable,colback=teal!4,colframe=teal!35!black,title={Before you move on}]
What does \sk{विपत्तिः} mean? (\textquotedblleft{}Adversity; a calamity\textquotedblright{}.) Say it once more.
\end{tcolorbox}
\section[\texorpdfstring{\sk{आपद्} (āpad)}{आपद् (āpad)}]{\sk{आपद्} (āpad) — a calamity}

\label{lesson:SA-C375-apad}



\begin{quote}
Before the new one: say the Sanskrit for ceases, then the Sanskrit for adversity.
\end{quote}

\subsection*{You'll want to know: \sk{आपद्}}
\textbf{\sk{आपद्}} — \emph{āpad} — \textquotedblleft{}a calamity\textquotedblright{}.

\begin{quote}
\sk{यदि} \sk{आपत्} \sk{आगच्छति}, \sk{तर्हि} \sk{धैर्यं} \sk{धारयतु।}
\end{quote}

\begin{quote}
\emph{yadi āpat āgacchati, tarhi dhairyaṃ dhārayatu.}
\end{quote}

\begin{quote}
If a calamity comes, keep your courage.
\end{quote}

A feminine noun.

The stem ends in \emph{-d}; said on its own, as here, the word ends in \emph{-t}: \sk{आपत्}, \emph{āpat}.

\begin{cousinweb}
\textbf{\sk{आपद्}} (\emph{āpad}) is built on \emph{ā} and the root √pad, \textquotedblleft{}to fall\textquotedblright{}; MW gives \emph{āpad} \textquotedblleft{}misfortune, calamity, distress\textquotedblright{} (MW).
\end{cousinweb}

\begin{grammarlens}[title={What you've built}]
One more word for saying what happens if.
\end{grammarlens}

\subsection*{Your turn}
\begin{itemize}
\raggedright
  \item \emph{Say it:} \emph{āpad}
  \item \emph{Say it:} \emph{āpad}, once more
  \item \emph{Say it:} \emph{yadi āpat āgacchati, tarhi dhairyaṃ dhārayatu.} — If a calamity comes, keep your courage.
  \item \emph{Recall:} say the Sanskrit for ceases, then the Sanskrit for adversity, then say \emph{āpad} again
\end{itemize}

\begin{tcolorbox}[breakable,colback=teal!4,colframe=teal!35!black,title={Before you move on}]
What does \sk{आपद्} mean? (\textquotedblleft{}A calamity\textquotedblright{}.) Say it once more.
\end{tcolorbox}
\section[\texorpdfstring{\sk{विनाशः} (vināśaḥ)}{विनाशः (vināśaḥ)}]{\sk{विनाशः} (vināśaḥ) — destruction; ruin}

\label{lesson:SA-C375-vinasah}



\begin{quote}
Before the new one: say the Sanskrit for adversity, then the Sanskrit for a calamity.
\end{quote}

\subsection*{You'll want to know: \sk{विनाशः}}
\textbf{\sk{विनाशः}} — \emph{vināśaḥ} — \textquotedblleft{}destruction; ruin\textquotedblright{}.

\begin{quote}
\sk{यदि} \sk{युद्धं} \sk{भवति}, \sk{तर्हि} \sk{विनाशः} \sk{भविष्यति।}
\end{quote}

\begin{quote}
\emph{yadi yuddhaṃ bhavati, tarhi vināśaḥ bhaviṣyati.}
\end{quote}

\begin{quote}
If there is a war, there will be destruction.
\end{quote}

A masculine noun.

\begin{cousinweb}
\textbf{\sk{विनाशः}} (\emph{vināśaḥ}) is built on \emph{vi} and the root √naś, \textquotedblleft{}to perish\textquotedblright{}; MW gives \emph{vi-nāśa} \textquotedblleft{}utter loss, annihilation, destruction, decay\textquotedblright{} (MW).
\end{cousinweb}

\begin{grammarlens}[title={What you've built}]
One more word for saying what happens if.
\end{grammarlens}

\subsection*{Your turn}
\begin{itemize}
\raggedright
  \item \emph{Say it:} \emph{vināśaḥ}
  \item \emph{Say it:} \emph{vināśaḥ}, once more
  \item \emph{Say it:} \emph{yadi yuddhaṃ bhavati, tarhi vināśaḥ bhaviṣyati.} — If there is a war, there will be destruction.
  \item \emph{Recall:} say the Sanskrit for adversity, then the Sanskrit for a calamity, then say \emph{vināśaḥ} again
\end{itemize}

\begin{tcolorbox}[breakable,colback=teal!4,colframe=teal!35!black,title={Before you move on}]
What does \sk{विनाशः} mean? (\textquotedblleft{}Destruction; ruin\textquotedblright{}.) Say it once more.
\end{tcolorbox}
\section[\texorpdfstring{\sk{अनावृष्टिः} (anāvṛṣṭiḥ)}{अनावृष्टिः (anāvṛṣṭiḥ)}]{\sk{अनावृष्टिः} (anāvṛṣṭiḥ) — a drought}

\label{lesson:SA-C375-anavrstih}



\begin{quote}
Before the new one: say the Sanskrit for a calamity, then the Sanskrit for destruction.
\end{quote}

\subsection*{You'll want to know: \sk{अनावृष्टिः}}
\textbf{\sk{अनावृष्टिः}} — \emph{anāvṛṣṭiḥ} — \textquotedblleft{}a drought\textquotedblright{}.

\begin{quote}
\sk{यदि} \sk{अनावृष्टिः} \sk{भवति}, \sk{तर्हि} \sk{कृषकः} \sk{दुःखी} \sk{भविष्यति।}
\end{quote}

\begin{quote}
\emph{yadi anāvṛṣṭiḥ bhavati, tarhi kṛṣakaḥ duḥkhī bhaviṣyati.}
\end{quote}

\begin{quote}
If there is a drought, the farmer will suffer.
\end{quote}

A feminine noun.

\begin{cousinweb}
\textbf{\sk{अनावृष्टिः}} (\emph{anāvṛṣṭiḥ}) is built from \emph{an-}, \textquotedblleft{}not\textquotedblright{}, and \emph{ā-vṛṣṭi}, \textquotedblleft{}rain\textquotedblright{}, from the root √vṛṣ, \textquotedblleft{}to rain\textquotedblright{}; MW gives \emph{an-āvṛṣṭi} \textquotedblleft{}want of rain, drought\textquotedblright{} (MW).
\end{cousinweb}

\begin{grammarlens}[title={What you've built}]
One more word for saying what happens if.
\end{grammarlens}

\subsection*{Your turn}
\begin{itemize}
\raggedright
  \item \emph{Say it:} \emph{anāvṛṣṭiḥ}
  \item \emph{Say it:} \emph{anāvṛṣṭiḥ}, once more
  \item \emph{Say it:} \emph{yadi anāvṛṣṭiḥ bhavati, tarhi kṛṣakaḥ duḥkhī bhaviṣyati.} — If there is a drought, the farmer will suffer.
  \item \emph{Recall:} say the Sanskrit for a calamity, then the Sanskrit for destruction, then say \emph{anāvṛṣṭiḥ} again
\end{itemize}

\begin{tcolorbox}[breakable,colback=teal!4,colframe=teal!35!black,title={Before you move on}]
What does \sk{अनावृष्टिः} mean? (\textquotedblleft{}A drought\textquotedblright{}.) Say it once more.
\end{tcolorbox}
\section[\texorpdfstring{\sk{प्रतिकरोति} (pratikaroti)}{प्रतिकरोति (pratikaroti)}]{\sk{प्रतिकरोति} (pratikaroti) — remedies; counteracts}

\label{lesson:SA-C375-pratikaroti}



\begin{quote}
Before the new one: say the Sanskrit for destruction, then the Sanskrit for a drought.
\end{quote}

\subsection*{You'll want to know: \sk{प्रतिकरोति}}
\textbf{\sk{प्रतिकरोति}} — \emph{pratikaroti} — \textquotedblleft{}remedies; counteracts\textquotedblright{}.

\begin{quote}
\sk{यदि} \sk{रोगः} \sk{आगच्छति}, \sk{तर्हि} \sk{वैद्यः} \sk{तं} \sk{प्रतिकरोति।}
\end{quote}

\begin{quote}
\emph{yadi rogaḥ āgacchati, tarhi vaidyaḥ taṃ pratikaroti.}
\end{quote}

\begin{quote}
If an illness comes, the doctor remedies it.
\end{quote}

Past (he …): \sk{सः} \sk{रोगं} \sk{प्रतिकृतवान्}, \emph{saḥ rogaṃ pratikṛtavān}.

\begin{cousinweb}
\textbf{\sk{प्रतिकरोति}} (\emph{pratikaroti}) is built on \emph{prati}, \textquotedblleft{}against\textquotedblright{}, and the root √kṛ, \textquotedblleft{}to do\textquotedblright{}; MW gives \emph{prati-√kṛ} \textquotedblleft{}to counteract, to resist\textquotedblright{} and \textquotedblleft{}to repay, to requite\textquotedblright{} (MW).
\end{cousinweb}

\begin{grammarlens}[title={What you've built}]
One more word for saying what happens if.
\end{grammarlens}

\subsection*{Your turn}
\begin{itemize}
\raggedright
  \item \emph{Say it:} \emph{pratikaroti}
  \item \emph{Say it:} \emph{pratikaroti}, once more
  \item \emph{Say it:} \emph{yadi rogaḥ āgacchati, tarhi vaidyaḥ taṃ pratikaroti.} — If an illness comes, the doctor remedies it.
  \item \emph{Recall:} say the Sanskrit for destruction, then the Sanskrit for a drought, then say \emph{pratikaroti} again
\end{itemize}

\begin{tcolorbox}[breakable,colback=teal!4,colframe=teal!35!black,title={Before you move on}]
What does \sk{प्रतिकरोति} mean? (\textquotedblleft{}Remedies; counteracts\textquotedblright{}.) Say it once more.
\end{tcolorbox}
